% from ega4-14-fr.tex
\oldpage[IV-3]{219}

Notons que si $A$ et $B$ vérifient les conditions de l'énoncé, $B$ est un anneau
de valuation discrète complet (Bourbaki, \emph{Alg. comm.}, chap.~III, \S~3,
n\textsuperscript{o}~3, prop.~7 et chap.~IV, \S~2, n\textsuperscript{o}~2,
cor.~3 de la prop.~9) de corps résiduel isomorphe à celui de $A$, donc
algébriquement clos, et puisque $u(z')=x$, il y a dans $X''=X\times_Y Z'$ un point
$x''$, dont les projections dans $X$ et $Z'$ sont $x$ et $z'$ et qui est rationnel
sur $k(z')$~; en outre, puisqu'il existe une $Z'$-section $v$ de $X''$ telle que
$v(z')=x''$, l'image par $v$ du point générique $s$ de $Z'$ est une générisation
$t$ de $x''$ dont la projection dans $Z'$ est $s$~; en appliquant (14.3.6), on voit
que la condition de l'énoncé est \emph{suffisante}. Prouvons maintenant qu'elle est
\emph{nécessaire}. Posons $X'=X\times_Y Y'$, $f'=f_{(Y')}:X'\to Y'$. Il y a par
hypothèse un point $x'\in X'$ au-dessus de $x$ et de $y'$ et rationnel sur $k(y')$
(\textbf{I}, 3.3.14), donc fermé dans $f'^{-1}(y')$. En vertu de (14.3.6), il y a
une composante irréductible $T'$ de $X'$ contenant $x'$ et dominant $Y'$, donc
((14.3.8) et (14.3.13)) la restriction $T'\to Y'$ de $f'$ est équidimensionnelle au
point $x$. On déduit donc de (14.5.5) qu'il y a une $A$-algèbre finie $B$ qui est
un anneau local intègre et intégralement clos et domine $A$ (donc est un anneau de
valuation discrète) et, en posant $Z'=\operatorname{Spec}(B)$ et
$X''=X\times_Y Z'=X'\times_{Y'}Z'$, une $Z'$-section $v:Z'\to X''$ telle que
l'image du point fermé $z'$ de $Z'$ par le morphisme composé $Z'\to X''\to X'$ soit
$x'$~; le morphisme composé $u:Z'\to X''\to X$ répond donc à la question.

\begin{proposition}[14.5.9]
\label{IV.14.5.9-fr}
Soient $Y$ un préschéma localement noethérien irréductible, $f:X\to Y$ un
morphisme localement de type fini, $y$ un point géométriquement unibranche de $Y$.
Alors les conditions suivantes sont équivalentes~:
\begin{enumerate}[label=\emph{\alph*)}]
  \item $f$ est universellement ouvert en tout point de $f^{-1}(y)$.
  \item $f$ est ouvert en tout point de $f^{-1}(y)$.
  \item Pour toute composante irréductible $Z$ de $f^{-1}(y)$, de point générique
  $z$, il existe une composante irréductible $X_i$ de $X$ contenant $z$ et
  équidimensionnelle sur $Y$ au point $z$.
  \item Pour tout point fermé $x$ de $f^{-1}(y)$, il existe une partie irréductible
  $X'$ de $X$, localement fermée dans $X$, contenant $x$, et telle que la
  restriction $X'\to Y$ de $f$ soit un morphisme quasi-fini et dominant.
\end{enumerate}
\end{proposition}

L'équivalence de $a)$, $b)$ et $c)$ a déjà été démontrée (14.4.2). Pour prouver
que $a)$ entraîne $d)$, notons qu'en vertu de (14.4.2), $a)$ entraîne qu'il existe
une composante irréductible $Z$ de $X$ contenant $x$ et équidimensionnelle sur $Y$
au point $x$~; l'existence de $X'$ provient alors de (14.5.3) appliqué à la
restriction $Z\to Y$ de $f$. Inversement, supposons $d)$ vérifiée~; en vertu du
critère de Chevalley (14.4.4), la restriction $X'\to Y$ de $f$ est un morphisme
universellement ouvert au point $x$, et a fortiori $f$ est ouvert au point $x$~;
$f$ est donc ouvert en tous les points fermés de $X_y=f^{-1}(y)$. Mais $X_y$ est
un $k(y)$-préschéma localement de type fini, donc un préschéma de Jacobson~;
l'ensemble des points fermés de $X_y$ est donc \emph{dense} dans $X_y$ (10.3.1),
et il résulte de (14.1.4) que $f$ est ouvert en \emph{tous} les points de $X_y$, ce
qui achève de prouver que $d)$ entraîne $b)$.

Le résultat suivant a été mis en évidence par D.~Mumford~:

\begin{proposition}[14.5.10]
\label{IV.14.5.10-fr}
Soient $Y$ un préschéma noethérien, $f:X\to Y$ un morphisme universellement ouvert,
surjectif et localement de type fini. Alors il existe un morphisme fini surjectif
$g:Y'\to Y$ tel que, si l'on pose $X'=X\times_Y Y'$ et
$f'=f_{(Y')}:X'\to Y'$, tout point $y'\in Y'$ admette un voisinage ouvert $U'$ tel
qu'il existe une $U'$-section de $f'^{-1}(U')$.
\end{proposition}

Nous démontrerons la proposition en plusieurs étapes.

